% ED26-REQ-N-d-023-03: Mac Lane--Moerdijk 1992, source-specific semantics.
\subsection{Mac Lane--Moerdijk: internal predicates and three forcing scopes}
\label{sec:p10-n-d-023-source-03}

Mac Lane and Moerdijk construct presheaf examples before interpreting typed
language and its generalized-element and site forcing semantics
\cite[I\S4, pp.~35--38; I\S6, pp.~44--47; VI\S5--7, pp.~296--318]
{MacLaneMoerdijk1992}. Here \(\mathcal C\) is small in the fixed universe,
\(\widehat{\mathcal C}=\mathsf{Set}^{\mathcal C^{\mathrm{op}}}\), and
\(x\cdot u=X(u)(x)\) for \(x\in X(c)\), \(u:d\to c\).
The source's \(\Rightarrow\) denotes Heyting implication, not material
implication; \(f^{-1}\) on subobjects is the host's \(f^*\).

The elementary-topos package is Definition~\ref{def:p10-topos}.
Its presheaf realization is supplied by Theorems~\ref{thm:p10-presheaf-limits},
\ref{thm:p10-presheaf-mono-epi}, \ref{thm:p10-presheaf-exponentials},
\ref{thm:p10-presheaf-topos} and \ref{thm:p10-presheaf-power-objects}, with
reciprocal proofs in Annex items~\ref{annex:p10.psh.1}--\ref{annex:p10.psh.5}.
These are the foundation interfaces below, not proofs of the following
semantic claims. The source writes \(Q^P(c)=\operatorname{Nat}(\mathrm yc\times
P,Q)\), with evaluation \((\theta,x)\mapsto\theta_c(1_c,x)\), and sends \(x\)
to the characteristic sieve \(\{u:x\cdot u\in A\}\). These are the same
host constructions with \(P,Q\) renamed \(X,Y\), not pointwise function sets.

\begin{hbtheorem}{thm:p10-mlm-presheaf-predicates}{proved=annex:p10.mlm.1}{Heyting predicates and all-arrow quantifiers}
For subpresheaves \(A,B\subseteq X\), top and bottom are \(X\) and \(0\),
and meets and joins are componentwise intersections and unions. Put
\[
 (A\Rightarrow B)(c)=\{x\in X(c):\forall u:d\to c,\
          x\cdot u\in A(d)\ \Longrightarrow\ x\cdot u\in B(d)\}.
\]
Then \(R\subseteq(A\Rightarrow B)\) iff \(R\cap A\subseteq B\).
Thus \(\operatorname{Sub}(X)\) is a Heyting algebra, with
\((\neg A)(c)=\{x:\forall u:d\to c,\ x\cdot u\notin A(d)\}\).
For any natural map \(f:X\to Y\), inverse image preserves these Heyting
operations and has both order adjoints \(\exists_f\dashv f^*\dashv\forall_f\):
\[
\begin{aligned}
 (\exists_f A)(c)&=\{y\in Y(c):\exists x\in A(c),\ f_c(x)=y\},\\
 (\forall_f A)(c)&=\{y\in Y(c):\forall u:d\to c\ \forall x\in X(d),\
                  f_d(x)=y\cdot u\ \Longrightarrow\ x\in A(d)\}.
\end{aligned}
\]
Frobenius says \(\exists_f(A\cap f^*B)=(\exists_f A)\cap B\).
For a pullback \(P=X\times_Y Z\) with projections \(p:P\to X\), \(q:P\to Z\),
and \(g:Z\to Y\), both Beck--Chevalley identities hold:
\[
             g^*\exists_f=\exists_qp^*,\qquad
             g^*\forall_f=\forall_qp^*.
\]
All formulas include empty components and the empty base.
\end{hbtheorem}

Implication and negation are I\S8 Proposition~5 (18)--(19), p.~56, where
further verification is left to the reader. The correctly typed universal
image is III\S8 (15), p.~148. Frobenius and Beck--Chevalley are IV\S9 (2),
Proposition~1 and (13), pp.~204--209 \cite{MacLaneMoerdijk1992}.
The annex supplies direct componentwise proofs; the source's general-topos
route instead uses cartesian-closed slices and stable epi--mono images.

\begin{hbdefinition}{def:p10-mlm-internal-operations}{Internal operations on power objects}
Give \(P(X)(c)=\operatorname{Sub}(\mathrm yc\times X)\) its Heyting structure.
Restrictions are pullback along \(\mathrm yu\times1_X\), so preserve that
structure. Thus component operations define natural internal maps
\(\wedge,\vee,\Rightarrow:P(X)\times P(X)\to P(X)\), \(\top,\bot:1\to P(X)\),
and \(\neg S=(S\Rightarrow\bot)\). Internal order is the equalizer of
\(\wedge\) and first projection: \(S\leq T\) means \(S\wedge T=S\).
For \(f:X\to Y\), define \(P(f):P(Y)\to P(X)\) and
\(\exists_f,\forall_f:P(X)\to P(Y)\), at component \(c\), using the external
inverse image and its adjoints along \(1_{\mathrm yc}\times f\).
On \(\operatorname{Nat}(Z,P(X))\cong\operatorname{Sub}(Z\times X)\), they act
by the corresponding external operations along \(1_Z\times f\).
This separates a map of power objects from an operation on an external
subobject poset. In particular \(P(1)\cong\Omega\) is an internal Heyting algebra.
\end{hbdefinition}

\begin{hbtheorem}{thm:p10-mlm-internal-representation}{proved=annex:p10.mlm.2}{Representation of the predicate operations}
These maps are natural, satisfy the Heyting identities and both internal
order adjunctions, and act on every parameter object exactly as stated.
Frobenius and both Beck--Chevalley identities are identities of natural maps,
not merely of global elements. This is the representability route of IV\S8
Theorem~1 (Internal), IV\S9 Theorem~2 and Proposition~4, pp.~201--209
\cite{MacLaneMoerdijk1992}.
\end{hbtheorem}

\begin{hbdefinition}{def:p10-mlm-language}{The Mitchell--B\'enabou language and interpretation}
Fix an elementary topos \(\mathcal E\). Its objects are types, including
\(1\), products, function types \(Y^X\), \(\Omega\), and \(P(X)=\Omega^X\).
A finite declared context \(\Gamma=(x_1:X_1,\ldots,x_n:X_n)\) denotes
\(G=\prod_iX_i\), with \(G=1\) if \(n=0\). A term of type \(Y\) denotes an
arrow \(G\to Y\); a variable denotes its projection. Keep contexts explicit,
including unused variables. In a common context the inductive clauses are:
\begin{itemize}
\item Pairing denotes \(\langle s,t\rangle\); projections and the unit term
      use their categorical arrows.
\item Each arrow \(h:X\to Y\) supplies \(h(s)\), denoting \(h\circ s\).
      In particular \(h:1\to Y\) supplies a closed term.
\item For \(r:G\to Y^X\), \(s:G\to X\), application \(r(s)\) denotes
      \(\operatorname{ev}\circ\langle r,s\rangle\).
\item For \(t:G\times X\to Y\), \(\lambda x.t\) denotes
      \(\Lambda t:G\to Y^X\); binding is capture-avoiding.
\item Equality of \(s,t:G\to X\) denotes
      \(\delta_X\circ\langle s,t\rangle\), where \(\delta_X\) classifies
      the diagonal \(X\to X\times X\).
\item For \(a:G\to X\), \(r:G\to P(X)\), membership \(a\in r\) denotes
      evaluation on \(\langle r,a\rangle\).
\item Formulas are \(\Omega\)-valued terms. Constants, connectives and
      negation use the internal Heyting maps on \(\Omega\).
\item For \(\phi\) in \(\Gamma,x:X\), its extension \([\phi]\hookrightarrow
      G\times X\) is the pullback of truth. Along \(\pi:G\times X\to G\),
      the extensions of \(\exists x\,\phi\), \(\forall x\,\phi\) are
      \(\exists_\pi[\phi]\), \(\forall_\pi[\phi]\). Equivalently transpose
      \(\phi:G\times X\to\Omega\) to \(G\to P(X)\), then apply internal
      quantifiers along \(X\to1\).
\item Comprehension \(\{x:X\mid\phi\}\) as a predicate-valued term denotes
      \(\Lambda\phi:G\to P(X)\). The same braces as an object denote
      \([\phi]\hookrightarrow G\times X\), not that arrow.
\end{itemize}
These are VI\S5 pp.~298--301 \cite{MacLaneMoerdijk1992}, expressed in one
declared context. The source merges ordered lists of free variables using
projections; repeated variables occur once, different variables of the same
type remain different factors. This context convention makes those
projections explicit without deleting unused possibly empty types.
Natural-number, recursion and arithmetic constructions require an NNO when
invoked; they are not consequences of the elementary-topos definition alone.
\end{hbdefinition}

\begin{hbtheorem}{thm:p10-mlm-contextual-semantics}{proved=annex:p10.mlm.3}{Substitution, comprehension and contextual validity in presheaves}
All language clauses are defined in \(\widehat{\mathcal C}\).
For capture-avoiding simultaneous typed substitution denoting \(a:H\to G\),
term interpretation is composition with \(a\), and \([\phi[a]]=a^*[\phi]\).
In \(\Gamma,x:X\), membership in the comprehension
\(\{x:X\mid\phi\}\) has exactly extension \([\phi]\).
Evaluation and transposition validate typed \(\beta\)- and \(\eta\)-identities.
Equality denotes the equalizer of the compared terms and supports replacing
equal terms throughout expressions. Predicate equality is extensional:
\(r=s\) has the same extension as
\(\forall x:X\,(x\in r\Leftrightarrow x\in s)\), with
\(\Leftrightarrow\) abbreviating the two Heyting implications.
Function equality is likewise equivalent to quantified pointwise equality.
For finite hypotheses \(\Delta\), semantic consequence in context \(G\)
means \(\bigcap_{\theta\in\Delta}[\theta]\subseteq[\psi]\); the empty
intersection is \(G\). Validity means \([\phi]=G\), or characteristic arrow
\(\mathsf{true}\,!_G\), and is equivalent to validity of universal closure
in context \(1\). No type inhabitation or choice principle is presumed.
\end{hbtheorem}

This theorem supplies semantic construction, not soundness of an unspecified
formal higher-order calculus. VI\S5 p.~302 sends specification of the proof
rules to Boileau--Joyal or Lambek--Scott \cite{MacLaneMoerdijk1992}.
The separately attributed Lambek--Scott1980 specification below supplies an
explicit dependent rule surface and soundness proof; it does not stand in
for the distinct 1986 book or Boileau--Joyal1981 presentation.

\begin{hbtheorem}{thm:p10-mlm-topos-structural-package}{cited=MacLaneMoerdijk1992}{Dependencies of generalized forcing}
In an elementary topos, subobject posets are Heyting algebras; pullback
preserves their Heyting structure and has both adjoints, with Beck--Chevalley.
Joins and existential images are images of coproduct and projection maps.
Epi--mono images and finite coproducts are stable under pullback; a monic epi
is an isomorphism; the initial object is strict. These are the earlier
structural results used by VI\S6: IV\S6--7, IV\S8 Theorem~1, and IV\S9.
They are not assumptions asserting the forcing clauses. For presheaves their
relevant constructions are provided by the five foundation proofs and
Theorem~\ref{thm:p10-mlm-presheaf-predicates}.
\end{hbtheorem}

\begin{hbdefinition}{def:p10-mlm-generalized-forcing}{Forcing at a generalized element}
For \(\alpha:U\to G\), \(U\forces\phi(\alpha)\) means that \(\alpha\) factors
through \([\phi]\hookrightarrow G\), equivalently
\(\operatorname{Im}(\alpha)\subseteq[\phi]\).
For a closed formula use \(U\to1\); for multiple variables use their tuple.
This is VI\S6 (1)--(2), p.~303 \cite{MacLaneMoerdijk1992}.
Here \(\forces\) is semantic, not derivability or set-theoretic forcing.
\end{hbdefinition}

\begin{hbtheorem}{thm:p10-mlm-generalized-clauses}{proved=annex:p10.mlm.4}{Generalized-element Kripke--Joyal clauses}
For \(\alpha:U\to G\), forcing persists under every \(p:V\to U\) and
descends along an epimorphism. Its clauses are:
\begin{itemize}
\item Top always holds; bottom holds iff \(U\cong0\).
\item Conjunction holds iff both conjuncts hold at \(\alpha\).
\item Disjunction holds iff there are \(p:V\to U\), \(q:W\to U\) with
      \([p,q]:V+W\to U\) epic, \(V\forces\phi(\alpha p)\) and
      \(W\forces\psi(\alpha q)\).
\item Implication holds iff every \(p:V\to U\) with
      \(V\forces\phi(\alpha p)\) also has \(V\forces\psi(\alpha p)\).
\item Negation holds iff every \(p:V\to U\) with
      \(V\forces\phi(\alpha p)\) has \(V\cong0\).
\item Existence \(U\forces\exists x:X\,\phi\) holds iff there are an epic
      \(p:V\to U\) and \(b:V\to X\) with \(V\forces\phi(\alpha p,b)\).
\item Universality \(U\forces\forall x:X\,\phi\) holds iff every \(p:V\to U\)
      and \(b:V\to X\) has \(V\forces\phi(\alpha p,b)\), equivalently
      \(U\times X\forces\phi(\alpha\pi_1,\pi_2)\).
\end{itemize}
Equality of terms \(s,t:G\to Y\) is forced iff \(s\alpha=t\alpha\).
Membership is forced iff the interpreted element--predicate tuple factors
through membership. These statements use the preceding structural package,
not enough points. They are VI\S6 Theorem~1 and Propositions~2--3,
pp.~304--310 \cite{MacLaneMoerdijk1992}; the annex expands the source's reader
verification for negation and membership too.
\end{hbtheorem}

\subsubsection{An explicit dependent rule specification: Lambek--Scott (1980)}
\label{sec:p10-ls1980-dependent-calculus}

Lambek and Scott's \(\mathcal L_1\) is a precise product/power type theory,
not literally the preceding exponential-term syntax. We use the actual
1980 article's \S1 pp.~216--218 and \S5 pp.~240--242
\cite{LambekScott1980ITT} as a separately qualified dependent specification.
The 1980 logical connectives are primitive and equality is Leibniz-defined;
neither the primitive-equality1983 version nor the 1986 book is merged into it.

\begin{hbdefinition}{def:p10-ls1980-language}{The source's product/power grammar and its declared fragments}
The full pure source language has basic types \(1,N,\Omega\), products
\(A\times B\), and power types \(PA\), with countably many variables of each
type. Its ten term clauses are: (1) a typed variable; (2) \(*:1\);
(3) \(\langle a,b\rangle:A\times B\) from \(a:A,b:B\);
(4) \(a\in u:\Omega\) from \(a:A,u:PA\);
(5) \(\{x:A\mid\phi(x)\}:PA\) from an \(\Omega\)-term, with binding of \(x\);
(6) \(0:N\); (7) \(Sn:N\) from \(n:N\);
(8) \(\top,\bot:\Omega\); (9) \(p\wedge q,p\vee q,p\Rightarrow q:\Omega\);
(10) \(\forall x:A\,\phi\), \(\exists x:A\,\phi:\Omega\).
Formulas are the terms of type \(\Omega\), and
\[
 \neg p=(p\Rightarrow\bot),\qquad
 p\Leftrightarrow q=(p\Rightarrow q)\wedge(q\Rightarrow p),\qquad
 a=_Ab=\forall u:PA\,(a\in u\Leftrightarrow b\in u),
\]
where the last \(u\) is fresh. There is no primitive equality term,
projection term, function type, lambda binder or application term in this
pure grammar. All binding and substitution use the source's free-for
condition, here made explicitly capture-avoiding.

The nonarithmetic fragment \(\mathcal L_1^-\) omits \(N,0,S\) and Peano
axioms; it is an explicitly identified restriction, not the full source
language renamed. The arithmetic extension retains exactly those source
constructors and axioms and requires an interpreted NNO. The optional
Boolean axiom is independent of that arithmetic extension.
The source's applied internal-language extension (\S5 p.~242) permits named
objects as types and named arrows as term operations: for an interpreted
\(h:A\to B\), \(h(a):B\) denotes composition. This extension allows arbitrary
object types, including uninhabited ones; it does not silently add a
primitive function/lambda grammar.
\end{hbdefinition}

\begin{hbdefinition}{def:p10-ls1980-contextual-validity}{Source-indexed entailment and interpretation}
The source relation \(p\vdash_Xq\) is indexed by a declared set \(X\) of
typed variables, containing both formulas' free variables; it is the least
relation satisfying the displayed structural, logical and nonlogical rows
of \ref{sheet:p10-internal-hol-ls1980}. The sheet splits each source logical
``if and only if'' into its rule directions and makes quantifier freshness
explicit. \(\vdash_Xp\) abbreviates \(\top\vdash_Xp\), not truth in a smaller
undeclared context.

Interpret types as the corresponding objects, products and power objects;
interpret the ten constructors by projections, pairing, membership,
comprehension, internal Heyting operations and quantification as above.
For possibly infinite \(X\), a generalized assignment at \(U\) is a family
\((a_x:U\to A_x)_{x\in X}\). Each term has finite support, so is evaluated
using the finite tuple for its free variables; no infinite context product
is presumed. Define \(p\models_Xq\) to mean: at every object \(U\) and every
such assignment, forcing of \(p\) implies forcing of \(q\).
For a finite \(X\), this is exactly
\([p]_G\subseteq[q]_G\) in its full declared context \(G=\prod_{x\in X}A_x\).
Generalized assignments retain unused declarations; if no assignment exists
at a noninitial \(U\), the corresponding test is vacuous, not an instruction
to delete its empty types.
\end{hbdefinition}

\begin{hbtheorem}{thm:p10-ls1980-leibniz-equality}{proved=annex:p10.mlm.7}{The Leibniz--diagonal interpretation bridge}
In an elementary topos, using the explicit structural package of
Theorem~\ref{thm:p10-mlm-topos-structural-package}, the source's Leibniz
formula \(a=_Ab\) has exactly the diagonal-equality extension for every pair
of contextual terms. The substitution and comprehension interpretations
therefore agree with the source's \S5 diagonal interpretation, pp.~241--242
\cite{LambekScott1980ITT}, without changing the source's defined syntax.
The same statement holds concretely in the already constructed presheaf
models, where the structural ingredients have the linked component proofs.
\end{hbtheorem}

% Lambek--Scott 1980 L_1, section 1 pp.216--218; biconditionals split into rules.
\calculussheet{sheet:p10-internal-hol-ls1980}{Lambek--Scott 1980 product/power type theory}{Source: \cite[\S1, Definitions~1.1--1.2 and axioms, pp.~216--218]{LambekScott1980ITT}. Types are \(1,N,\Omega\), closed under \(A\times B\) and \(PA\); terms are exactly the ten clauses in Definition~\ref{def:p10-ls1980-language}. The explicitly nonarithmetic fragment omits \(N,0,S\) and Peano rows; the applied extension names additional interpreted object types and arrows as in source \S5 p.~242. Judgments \(p\vdash_Xq\) have formula free variables contained in the declared set \(X\); \(\vdash_Xp\) abbreviates \(\top\vdash_Xp\). All substitutions are typed and free-for/capture-avoiding. Write \(\neg p=(p\Rightarrow\bot)\), \(p\Leftrightarrow q=(p\Rightarrow q)\wedge(q\Rightarrow p)\), and \(a=_Ab\) for \(\forall u:PA\,(a\in u\Leftrightarrow b\in u)\), with \(u\) fresh. Equality is defined, not primitive; neither function/lambda types nor projections are primitive terms of pure \(\mathcal L_1\). Each source logical iff is displayed as its forward and reverse rules. Quantifier freshness is explicit; optional Peano rows require the NNO interpretation and the Boolean row requires Booleanness.}
\hbrule{rule:p10-internal-hol-ls1980-refl}{Reflexivity}{}{p\vdash_Xp}{\none}
\hbrule{rule:p10-internal-hol-ls1980-trans}{Transitivity}{p\vdash_Xq\qquad q\vdash_Xr}{p\vdash_Xr}{\none}
\hbrule{rule:p10-internal-hol-ls1980-weaken}{Variable weakening}{p\vdash_Xq}{p\vdash_Yq}{\(X\subseteq Y\); declared types are unchanged. No unused-variable deletion rule.}
\hbrule{rule:p10-internal-hol-ls1980-subst}{Typed substitution}{\phi(y)\vdash_{X\cup\{y:B\}}\psi(y)}{\phi(b)\vdash_X\psi(b)}{\(b:B\), \(\mathrm{FV}(b)\subseteq X\); \(b\) is free for \(y\) in both formulas. If \(y\in X\), its original assigned value is used when evaluating \(b\), and the premise receives the substituted value at \(y\).}
\hbrule{rule:p10-internal-hol-ls1980-top}{Logical (1)}{}{p\vdash_X\top}{\none}
\hbrule{rule:p10-internal-hol-ls1980-bottom}{Logical (2)}{}{\bot\vdash_Xp}{\none}
\hbrule{rule:p10-internal-hol-ls1980-meet-intro}{Logical (3), reverse}{p\vdash_Xq\qquad p\vdash_Xr}{p\vdash_Xq\wedge r}{\none}
\hbrule{rule:p10-internal-hol-ls1980-meet-left}{Logical (3), forward 1}{p\vdash_Xq\wedge r}{p\vdash_Xq}{\none}
\hbrule{rule:p10-internal-hol-ls1980-meet-right}{Logical (3), forward 2}{p\vdash_Xq\wedge r}{p\vdash_Xr}{\none}
\hbrule{rule:p10-internal-hol-ls1980-join-elim}{Logical (4), reverse}{p\vdash_Xr\qquad q\vdash_Xr}{p\vee q\vdash_Xr}{\none}
\hbrule{rule:p10-internal-hol-ls1980-join-left}{Logical (4), forward 1}{p\vee q\vdash_Xr}{p\vdash_Xr}{\none}
\hbrule{rule:p10-internal-hol-ls1980-join-right}{Logical (4), forward 2}{p\vee q\vdash_Xr}{q\vdash_Xr}{\none}
\hbrule{rule:p10-internal-hol-ls1980-residual-intro}{Logical (5), forward}{p\wedge q\vdash_Xr}{p\vdash_Xq\Rightarrow r}{\none}
\hbrule{rule:p10-internal-hol-ls1980-residual-elim}{Logical (5), reverse}{p\vdash_Xq\Rightarrow r}{p\wedge q\vdash_Xr}{\none}
\hbrule{rule:p10-internal-hol-ls1980-forall-intro}{Logical (6), reverse}{p\vdash_{X\cup\{y:B\}}\psi(y)}{p\vdash_X\forall y:B\,\psi(y)}{\(y\notin\mathrm{FV}(p)\); other free variables of \(\psi\) lie in \(X\). If \(y\in X\), its declared value is overridden when testing the bound variable.}
\hbrule{rule:p10-internal-hol-ls1980-forall-elim}{Logical (6), forward}{p\vdash_X\forall y:B\,\psi(y)}{p\vdash_{X\cup\{y:B\}}\psi(y)}{\(y\notin\mathrm{FV}(p)\); other free variables of \(\psi\) lie in \(X\).}
\hbrule{rule:p10-internal-hol-ls1980-exists-elim}{Logical (7), reverse}{\psi(y)\vdash_{X\cup\{y:B\}}p}{\exists y:B\,\psi(y)\vdash_Xp}{\(y\notin\mathrm{FV}(p)\); other free variables of \(\psi\) lie in \(X\).}
\hbrule{rule:p10-internal-hol-ls1980-exists-intro}{Logical (7), forward}{\exists y:B\,\psi(y)\vdash_Xp}{\psi(y)\vdash_{X\cup\{y:B\}}p}{\(y\notin\mathrm{FV}(p)\); other free variables of \(\psi\) lie in \(X\).}
\hbrule{rule:p10-internal-hol-ls1980-comprehension}{Nonlogical (a)}{}{\vdash_X\forall x:A\,(x\in\{x:A\mid\phi(x)\}\Leftrightarrow\phi(x))}{Parameters of \(\phi\) other than its displayed \(x\) lie in \(X\); binders are renamed to avoid capture, not to delete declared variables.}
\hbrule{rule:p10-internal-hol-ls1980-power-ext}{Nonlogical (b), power}{}{\vdash_{\varnothing}\forall u:PA\,\forall v:PA\,((\forall x:A\,(x\in u\Leftrightarrow x\in v))\Rightarrow u=_{PA}v)}{\(u,v,x\) are distinct bound variables.}
\hbrule{rule:p10-internal-hol-ls1980-proposition-ext}{Nonlogical (b), propositions}{}{\vdash_{\varnothing}\forall s:\Omega\,\forall t:\Omega\,((s\Leftrightarrow t)\Rightarrow s=_{\Omega}t)}{\(s,t\) are distinct bound variables.}
\hbrule{rule:p10-internal-hol-ls1980-unit}{Nonlogical (c), unit}{}{\vdash_{\varnothing}\forall x:1\,x=_1*}{\none}
\hbrule{rule:p10-internal-hol-ls1980-product-surj}{Nonlogical (c), decomposition}{}{\vdash_{\varnothing}\forall z:A\times B\,\exists x:A\,\exists y:B\,z=_{A\times B}\langle x,y\rangle}{\(z,x,y\) are distinct bound variables.}
\hbrule{rule:p10-internal-hol-ls1980-product-inj}{Nonlogical (c), injectivity}{}{\vdash_{\varnothing}\forall x:A\,\forall x':A\,\forall y:B\,\forall y':B\,(\langle x,y\rangle=_{A\times B}\langle x',y'\rangle\Rightarrow(x=_Ax'\wedge y=_By'))}{All four bound variables are distinct.}
\hbrule{rule:p10-internal-hol-ls1980-peano-zero}{Optional nonlogical (d), zero}{}{\vdash_{\varnothing}\forall x:N\,(Sx=_N0\Rightarrow\bot)}{Arithmetic extension only; \(N\) is interpreted as an NNO with its zero and successor.}
\hbrule{rule:p10-internal-hol-ls1980-peano-inj}{Optional nonlogical (d), successor}{}{\vdash_{\varnothing}\forall x:N\,\forall y:N\,(Sx=_NSy\Rightarrow x=_Ny)}{Arithmetic extension only; same NNO interpretation.}
\hbrule{rule:p10-internal-hol-ls1980-peano-induction}{Optional nonlogical (d), induction}{}{\vdash_{\varnothing}\forall u:PN\,((0\in u\wedge\forall x:N\,(x\in u\Rightarrow Sx\in u))\Rightarrow\forall y:N\,y\in u)}{Arithmetic extension only; same NNO interpretation. Bound \(u,x,y\) are distinct.}
\hbrule{rule:p10-internal-hol-ls1980-boolean}{Optional Boolean axiom}{}{\vdash_{\varnothing}\forall t:\Omega\,(t\vee\neg t)}{Classical extension only; the interpreting topos is Boolean. Independent of whether Peano rows are included.}


\begin{hbtheorem}{thm:p10-ls1980-hol-soundness}{proved=annex:p10.mlm.8}{Rule-by-rule intuitionistic higher-order soundness}
Every derivation using the nonarithmetic intuitionistic rows (excluding both
optional extensions) of
\ref{sheet:p10-internal-hol-ls1980}, in either the pure or the declared
object/arrow-named applied language, satisfies
\[
                 p\vdash_Xq\quad\Longrightarrow\quad p\models_Xq.
\]
In finite contexts this is the subobject inclusion of
Definition~\ref{def:p10-ls1980-contextual-validity}; in sheaf models it uses
the site forcing clauses, not pointwise Boolean satisfaction. The proof
validates every rule direction and nonlogical axiom before induction on the
derivation. An arbitrary elementary-topos interpretation uses only the
identified earlier structural package, not soundness as an assumed premise.

The full arithmetic rows are also sound when \(N,0,S\) are interpreted as an
NNO and its zero and successor; the annex derives the three Peano properties
from that universal property. Presheaves have the explicitly constructed
constant natural-number NNO. The Boolean row is sound precisely under
Booleanness, by Theorem~\ref{thm:p10-mlm-internal-epi-boolean-choice}.
Neither NNO nor Booleanness is required for the nonarithmetic intuitionistic
fragment. No completeness, normalization, realizability or free-topos
metatheorem is claimed here.
\end{hbtheorem}

\begin{hbtheorem}{thm:p10-ls1980-exponential-graph-bridge}{proved=annex:p10.mlm.9}{The graph bridge to exponential terms}
Inside \(P(A\times B)\), let \(\operatorname{Fun}(A,B)\) be the subobject
defined, in the source's product/power language, by
\[
 \forall x:A\,\exists y:B\left(\langle x,y\rangle\in R\wedge
       \forall y':B\,(\langle x,y'\rangle\in R\Rightarrow y'=_{B}y)\right).
\]
There is a canonical isomorphism \(\operatorname{Fun}(A,B)\cong B^A\),
natural on all parameter objects. A function is sent to its graph, and a
total single-valued relation is sent to its unique arrow. Under it membership
in the graph is equality with evaluation. A lambda term's graph is the
source comprehension
\[
 \{z:A\times B\mid\exists x:A\,\exists y:B\,
             (z=_{A\times B}\langle x,y\rangle\wedge y=_Bt(x))\}.
\]
It represents exactly the exponential transpose of the interpreted body.
Graph extensionality and the inverse representation give the typed beta and
eta conversions, including empty types. Thus the larger exponential-term
interpretation has an actual object/arrow bridge to the applied source
language. This is not a claim that function types or lambdas were primitive
in \(\mathcal L_1\), nor a proof of syntactic conservativity of a pure-language
extension. If exponential-term macros or their beta/eta conversions are used
with the source logical rules, their denotation-preserving expansion and the
same rule validation preserve soundness.
\end{hbtheorem}

Freshness is consequential. Allowing \(y\) to remain a parameter of \(p\)
in logical (6) would infer \(p\vdash_X\forall y\,p\) from reflexivity when
\(p\) depends on \(y\); in a two-element set with \(p(y)\) selecting one
element that conclusion is false. Conversely deleting an unused variable
of empty type changes the contextual meaning. Neither maneuver is a rule
of the declared interpretation.

\begin{hbdefinition}{def:p10-mlm-site-setup}{Site setup and its separate construction boundary}
A Grothendieck topology \(J\) assigns covering sieves \(J(c)\), contains the
maximal sieve, is stable under inverse image, and is transitive: if \(S\in J(c)\)
and \(u^*R\in J(d)\) for every \(u:d\to c\) in \(S\), then \(R\in J(c)\).
A sheaf uniquely amalgamates matching families on covers. A subsheaf has
local membership: a section belongs iff its restrictions on a cover belong.
Take the already established \(\mathcal E=\operatorname{Sh}(\mathcal C,J)\)
and associated-sheaf adjunction \(a\dashv i\); their construction is outside
this overview. Yoneda and the adjunction give
\[
 X(c)\cong\operatorname{Nat}(\mathrm yc,iX)\cong\mathcal E(a(\mathrm yc),X).
\]
For \(\gamma\in G(c)\), define \(c\forces\phi(\gamma)\) by
\(\gamma\in[\phi](c)\). It is generalized forcing from \(a(\mathrm yc)\),
not from a possibly nonsheaf \(\mathrm yc\). This identification is VI\S7
(1)--(3), pp.~315--316 \cite{MacLaneMoerdijk1992}; no sheafification or
geometric-morphism construction is asserted here.
\end{hbdefinition}

\begin{hbtheorem}{thm:p10-mlm-site-clauses}{proved=annex:p10.mlm.5}{Site and presheaf stage clauses}
At a site stage, equality and membership use their interpreted components,
conjunction means both at \(c\), top always holds, and bottom holds iff
\(\varnothing\in J(c)\). For \(\gamma\in G(c)\):
\begin{itemize}
\item Disjunction holds iff some covering sieve \(S\) has, for every
      \(u:d\to c\) in \(S\), either \(d\forces\phi(\gamma\cdot u)\) or
      \(d\forces\psi(\gamma\cdot u)\).
\item Implication holds iff every \(u:d\to c\) forcing the antecedent at
      \(\gamma\cdot u\) also forces the consequent there.
\item Negation holds iff every positive-forcing \(u:d\to c\) has
      \(\varnothing\in J(d)\).
\item Existence holds iff some covering sieve \(S\) has for each \(u:d\to c\)
      in \(S\) a \(b\in X(d)\) with \(d\forces\phi(\gamma\cdot u,b)\).
\item Universality holds iff for every \(u:d\to c\) and \(b\in X(d)\),
      \(d\forces\phi(\gamma\cdot u,b)\).
\end{itemize}
Forcing persists under restriction and is local on covers. At the trivial
topology, only the maximal sieve covers: disjunction then has a disjunct at
\(c\), existence a witness in \(X(c)\), and negation excludes every
positive-forcing restriction. Bottom is never forced at a base object.
Implication and universality still test all arrows, including identity.
No products in \(\mathcal C\) are required.
\end{hbtheorem}

These are VI\S7 Theorem~1 and (ii\(^{\prime}\)), (iv\(^{\prime}\)),
(v\(^{\prime}\)), pp.~316--318 \cite{MacLaneMoerdijk1992}.
The source's first proof route uses explicit subsheaf operations from III\S8;
the annex expands it. The other uses generalized forcing and the generation
and cover-detection properties of sheafified representables (III\S6,
III\S7 Corollary~7). Those are separate dependencies, not an enough-points
assumption. The generic-element test \(U\times X\) does not become one base
stage without further hypotheses; VI\S7 Remark~2 assumes products in
\(\mathcal C\) and a quantified type \(a(\mathrm yB)\).

\begin{hbexample}{def:p10-mlm-arrow-failures}{Two stages and the false pointwise substitute}
Let \(a:0\to1\) be the only nonidentity arrow of \(\mathcal C\).
In the terminal presheaf take \(A(0)=\{*\}\), \(A(1)=\varnothing\).
Then \(\neg A=0\), since restriction along \(a\) reaches \(A(0)\), and
\(\neg\neg A=1\neq A\). At \(1\), neither \(A\) nor \(\neg A\) holds.
Pointwise complement would put \(*\) at \(1\) but not \(0\), so is not a
subpresheaf. For \(f:A\to1\), \(\exists_f A=A\) as a subobject of \(1\),
whereas \(\forall_f0=0\). A component-only universal formula would falsely
put \(*\) at \(1\), because \(A(1)\) is empty; \(a\) and \(*\in A(0)\)
refute it. The classifier has two sieves at \(0\) and three at \(1\):
\(\varnothing\), \(\{a\}\), and \(\{a,1_1\}\). The middle sieve becomes
true after restriction but is not true at \(1\).
This matches the source's arrow-category example after reversing its arrow
for the host's contravariant convention \cite[I\S4, pp.~35--36; I\S8, p.~56]
{MacLaneMoerdijk1992}.
\end{hbexample}

\begin{hbexample}{def:p10-mlm-locality-failures}{Locality, choice, and unused types}
In \(\mathsf{Set}\), let \(U=\{0,1\}\), with predicates supported on the
respective singletons. Their disjunction holds at generalized element
\(1_U\), although neither disjunct does: the two inclusions have epic
copairing. The one presheaf base-stage representable for \(\mathsf{Set}\)
is a singleton, not \(U\), so its disjunction clause is different.
For a two-element group \(H\), the map from its regular right action to
the trivial singleton action is epic but has no equivariant section:
such a section would select a point fixed by the nonidentity element.
Internal existence is therefore not natural choice. Invariant subsets
nevertheless have invariant complements, so the subobject Heyting algebras
of \(H\)-sets are Boolean, as in I\S8 p.~56 \cite{MacLaneMoerdijk1992}.
Finally \(\bot\) is valid in declared context \(x:0\), whose context is
initial, but not in the empty context of a nondegenerate topos. Deleting
unused variables is not validity-preserving without inhabitation.
If the empty sieve covers a site object \(c\), every sheaf has one section
there and bottom is forced; bottom is not universally absent at site stages.
\end{hbexample}

\begin{hbtheorem}{thm:p10-mlm-internal-epi-boolean-choice}{proved=annex:p10.mlm.6}{Three internal-to-external translations}
For \(r:U\to Y^X\) with transpose \(\widetilde r:U\times X\to Y\),
\[
 U\forces\forall y:Y\,\exists x:X\,r(x)=y
 \ \Longleftrightarrow\
 \langle\pi_1,\widetilde r\rangle:U\times X\to U\times Y\text{ is epic}.
\]
A topos is Boolean iff \(\forall p:\Omega\,(p\vee\neg p)\) is valid.
For all types \(X,Y\), validity of
\[
 \forall f:Y^X\left((\forall y:Y\,\exists x:X\,f(x)=y)
       \Rightarrow\exists s:X^Y\,\forall y:Y\,f(s(y))=y\right)
\]
is equivalent to every \((-)^E\) preserving epimorphisms, for every object
\(E\). This is the source's internal axiom of choice (IAC), not the external
assertion that every epi has a global section, and is not assumed here.
The translations are VI\S5 (7)--(10), VI\S6 (10)--(26), pp.~301--315
\cite{MacLaneMoerdijk1992}.
\end{hbtheorem}

For the dense topology on a poset, a sieve covers \(p\) iff every \(q\leq p\)
has an \(r\leq q\) in it. Thus disjunction and existence require, after
each refinement, a further refinement with a disjunct or witness.
Implication and universality test all \(q\leq p\); negation allows no
positive-forcing refinement. This is VI\S7's dense/double-negation example
on p.~318, not a proof of continuum-hypothesis independence or a transfer
back to set theory.
